\section{Introduction}
\label{sec:intro}

The public ALMA archive holds thousands of hours of millimetre and
submillimetre observations taken for other reasons, mostly studies of dust
discs and planet formation. Their visibilities allow a spectrum at the
position of the star at the centre of each field, at a resolution and depth
set by the original observer's correlator setup. This paper searches those
spectra for a narrowband transmitter at the star's own position. No such
signal was found.

Searches for \emph{technosignatures} --- detectable imprints of technology on
the electromagnetic spectrum, intended or not
\citep{Sheikh2020,SETIreview2025} --- have examined many thousands of stars
over six decades, yet the last published accounting of that effort covers
only \HaystackFrac{} of the space of transmitter position, power, frequency,
bandwidth, polarisation and repetition rate
\citep{Wright2018}, and essentially all of it lies below
${\sim}30$\,GHz. Two targeted searches have reached nearby stars above that frequency.
\citet{Steffes1994} observed \SbPriorNStarA{} solar-type stars within
\SbPriorDistA\,pc at \SbPriorFreqGHz\,GHz, and \citet{Mauersberger1996}
included $\tau$~Cet and $\epsilon$~Eri among \SbPriorNStarB{} targets at the
same frequency; both searched a single line about \SbPriorBandMHz\,MHz wide at
one epoch, and \citet{Mason2025} has since searched more distant stars in
Band~3. This survey covers
\DnuA\,GHz in \CovNIslands{} disjoint intervals, searches a grid of drift rates rather
than a fixed frequency, tests every crossing against the \BdNCtrl{} control positions of
its own field, and uses repeat observations to test persistence.

The band has physical advantages and physical costs. Interstellar dispersion
and scintillation are negligible at these frequencies, so a carrier arrives
essentially undistorted, and a given aperture delivers far more gain per watt
supplied. Against that, receiver system temperatures are higher, the
atmosphere is opaque in places at both ends of the link, and the same high
gain narrows the beam: the \BenchDiam\,m benchmark of
\S\ref{sec:excludes} radiates into a beam only \BmArcsec{} arcsec across, so
the Earth must lie inside a beam that was deliberately aimed. What else makes
the band hard belongs to the archive rather than to the sky, and
\S\ref{sec:background} sets that out.

The signal sought is an \emph{unresolved spectral carrier} --- emission confined
to a few channels, unresolved at the stellar position, and present at more
than one epoch. Recurrence rather than brightness is what separates an
interesting event from a detection \citep{Sheikh2021}, and the archive holds
a second, separated epoch for some of these systems and not for others, so
what could be searched and what could be confirmed are different extents. The
frequency is allowed to drift, as it must for a transmitter carried on a
rotating, orbiting body (\S\ref{sec:statistic}). ``Narrowband'' describes the
assumed transmitter and not the resolving power of the archived channels. The
targets were chosen by other programmes for other purposes, so every
statistical statement here is about an archive-selected sample
(\S\ref{sec:bias}).
